% =====================================================================
% Bitácora de desarrollo de Franui · Polykat
% Registro de aprendizajes, decisiones, evidencias y avances del proyecto.
%
% Compilar (desde esta carpeta, dos veces para el índice):
%   xelatex bitacora.tex && xelatex bitacora.tex
% Entradas nuevas: entradas/AAAA-MM-DD-tema.tex, y agregar su \input al final.
% =====================================================================
\documentclass[11pt,a4paper]{article}

\usepackage{fontspec}
\setmainfont{Liberation Sans}           % mismas medidas que Arial
\usepackage[spanish,es-noshorthands,es-tabla]{babel}
\usepackage[a4paper,margin=2.2cm]{geometry}
\usepackage{xcolor}
\definecolor{carmesi}{HTML}{9E2A3A}      % paleta Higo de la app
\definecolor{tinta}{HTML}{2A1715}
\definecolor{suave}{HTML}{6B4B45}
\definecolor{rosa}{HTML}{F7EBE8}
\color{tinta}

\usepackage{titlesec}
\titleformat{\section}{\LARGE\bfseries\color{carmesi}}{\thesection}{.6em}{}
\titleformat{\subsection}{\large\bfseries}{\thesubsection}{.6em}{}
\titleformat{\subsubsection}{\normalsize\bfseries\color{suave}}{}{0em}{}
\titlespacing*{\section}{0pt}{0pt}{10pt}

\usepackage{enumitem}
\setlist{itemsep=2pt,topsep=4pt}
\usepackage{booktabs,tabularx,array}
\renewcommand{\arraystretch}{1.25}
\usepackage[most]{tcolorbox}
\newtcolorbox{idea}[1][]{colback=rosa,colframe=carmesi,boxrule=0pt,leftrule=3pt,arc=2pt,left=8pt,right=8pt,top=6pt,bottom=6pt,#1}
\newtcolorbox{hipotesis}[1][Hipótesis por validar]{colback=white,colframe=suave,boxrule=.6pt,arc=3pt,
  title={#1},fonttitle=\bfseries\small,coltitle=white,colbacktitle=suave,left=8pt,right=8pt}
\newtcolorbox{formulacion}{colback=white,colframe=carmesi,boxrule=.8pt,arc=3pt,left=10pt,right=10pt,top=8pt,bottom=8pt}

\usepackage{fancyhdr}
\pagestyle{fancy}
\fancyhf{}
\renewcommand{\headrulewidth}{0pt}
\fancyfoot[L]{\small\color{suave}Franui · Polykat — Bitácora de desarrollo}
\fancyfoot[R]{\small\color{suave}\thepage}

\usepackage[hidelinks]{hyperref}
\usepackage{csquotes}

% --- Comandos de la bitácora ---
% \entrada{Título}{Fecha}{Tipo: Clase | Avance | Evidencia | Decisión}
\newcommand{\entrada}[3]{\clearpage\section{#1}\noindent{\small\color{suave}\textbf{#2}\quad·\quad #3}\par\medskip}
% Dato que falta: queda marcado en carmesí hasta que lo investiguemos con una fuente
\newcommand{\porinvestigar}[1]{\textcolor{carmesi}{\textbf{[por investigar: #1]}}}
% Conector de la sintaxis del problema (en carmesí, como en clase)
\newcommand{\conector}[1]{\textcolor{carmesi}{\textbf{#1}}}

\begin{document}

% ------------------------------- Portada -------------------------------
\begin{titlepage}
  \vspace*{4cm}
  {\color{suave}\large\itshape Polykat}\par
  {\fontsize{46}{50}\selectfont\bfseries F R A N U I}\par
  \vspace{2mm}{\color{carmesi}\rule{8cm}{1pt}}\par
  \vspace{6mm}{\Large Bitácora de desarrollo}\par
  \vspace{2mm}{\color{suave} Aprendizajes, decisiones, evidencias y avances del proyecto}\par
  \vfill
  {\large Francisca Ignacia Sanhueza Flores}\par
  {\color{suave} Creadora de Franui · Polykat}\par
  \vspace{4mm}
  {\color{suave} Inicio de la bitácora: 10 de noviembre de 2026}\par
  {\color{suave} \url{https://polykat-franui.github.io}}
\end{titlepage}

\tableofcontents

% --------------------------- Cómo la usamos ---------------------------
\clearpage
\section*{Cómo usamos esta bitácora}
\addcontentsline{toc}{section}{Cómo usamos esta bitácora}

Franui es un proyecto personal. Esta bitácora no es una tarea: es el registro honesto de cómo
construimos la idea hasta convertirla en una empresa. Nos sirve para no perder lo aprendido, para
tomar mejores decisiones y, más adelante, como respaldo al postular a fondos concursables.

\subsection*{Tipos de entrada}
\begin{description}[leftmargin=2.6cm,style=nextline]
  \item[Clase] Conceptos de una clase o lectura y cómo se aplican a Franui.
  \item[Avance] Lo que cambió en la app o en la web, con su porqué.
  \item[Evidencia] Observaciones, entrevistas, encuestas o datos, y sus hallazgos.
  \item[Decisión] Una decisión importante, las opciones que había y por qué elegimos una.
\end{description}

\subsection*{Reglas}
\begin{itemize}
  \item Separar siempre \textbf{hechos} (lo que observamos o medimos) de \textbf{hipótesis} (lo que creemos).
  \item Ningún dato inventado: lo que falta queda marcado como \porinvestigar{qué falta}.
  \item Toda cifra externa lleva su fuente (APA 7).
  \item Las personas entrevistadas o que opinan aparecen sin nombre ni datos personales.
\end{itemize}

% --------------------------- Línea base ---------------------------
\clearpage
\section{Línea base: dónde está Franui hoy}
\noindent{\small\color{suave}\textbf{10 de noviembre de 2026}\quad·\quad Avance}\par\medskip

\begin{idea}
\textbf{Qué es.} Una app web instalable para aprender el vocabulario y las frases cotidianas de varios
idiomas a la vez, deslizando tarjetas en sesiones cortas, con repetición espaciada y una compañera,
Kitty, que acompaña y motiva.
\end{idea}

\subsection*{Producto (hechos)}
\begin{itemize}
  \item App y página web publicadas en \url{https://polykat-franui.github.io}.
  \item Banco de 4.996 palabras y frases en inglés, traducidas a español, portugués, francés, italiano
        y alemán (traducciones en revisión).
  \item Oraciones de ejemplo para las 1.000 palabras más usadas, en los 6 idiomas.
  \item Método: tarjetas que se deslizan, repetición espaciada, meta diaria de 20 palabras, racha,
        \enquote{Prueba de Kitty} al final de la sesión.
  \item Motivación: Kitty Tokens, diamantes, jardín de sakura que crece al estudiar, ligas semanales y ranking.
  \item Cuentas con Google y avance guardado en la nube (Firebase).
  \item Instalable en el celular; funciona sin conexión.
\end{itemize}

\subsection*{Primeras personas usuarias (hechos)}
\begin{itemize}
  \item 13 cuentas creadas con Google entre el 13 y el 14 de octubre de 2026 (registro de Firebase).
  \item Comentarios recibidos de compañeros por mensajería (ver entrada de la Clase 1 y 2, hallazgos).
\end{itemize}

\subsection*{Lo que todavía no sabemos}
\begin{itemize}
  \item Cuántas personas vuelven al día siguiente, a la semana y al mes (la app aún no lo mide).
  \item Quién pagaría por Franui y cuánto.
  \item Si el problema que creemos resolver es el que de verdad tienen las personas: por eso existe esta bitácora.
\end{itemize}

% ------------------------------ Entradas ------------------------------
% Entrada: Clase 1 y 2 del curso Proyectos (CD1201) — identificación y formulación de problemas.
\entrada{Clase 1 y 2: formular bien el problema}{10 de noviembre de 2026}{Clase}

\begin{idea}
\textbf{Idea central.} Antes de construir una solución hay que entender y formular bien el problema.
Los errores en esta etapa son los más caros (el puente Cau Cau y el espejo del Hubble son ejemplos).
Para Franui esto significa: ya tenemos un prototipo funcionando, pero aún no hemos validado el problema
con personas fuera de nuestro círculo cercano.
\end{idea}

\subsection{Conceptos clave}

\subsubsection{Desafío, problemática y problema}
\begin{itemize}
  \item \textbf{Desafío:} una oportunidad para aplicar capacidades; no implica algo negativo.
  \item \textbf{Problemática:} un conjunto de problemas relacionados que necesita un análisis amplio.
  \item \textbf{Problema:} una dificultad específica y concreta, que pide una solución directa.
\end{itemize}

\subsubsection{Qué nos impide formular bien un problema}
Información incompleta, sesgos (por ejemplo, el de confirmación), complejidad, apuro por encontrar
una solución y falta de experiencia. Para contrarrestarlos: pensamiento crítico, investigación,
análisis de datos, comunicación, pensamiento sistémico y paciencia.

\subsubsection{Elementos de un problema}
\begin{description}[leftmargin=2.8cm,style=nextline]
  \item[Usuario y contexto] Quién tiene el problema. Se acota con segmentación: edad, ocupación,
        lugar, nivel socioeconómico, etc.
  \item[Necesidad] Lo que la persona requiere o le falta para alcanzar un objetivo.
  \item[Síntomas] Señales visibles de que algo no funciona; no explican la causa.
  \item[Evidencias] Datos objetivos y verificables que demuestran que el problema existe.
  \item[Hallazgos] Datos concretos que salen de nuestra propia investigación (\enquote{esto es lo que encontramos}).
  \item[Insights] La interpretación profunda de los hallazgos: por qué ocurre y dónde está la oportunidad.
\end{description}

\subsubsection{Matriz Aplicada del Problema (MAP) y sintaxis}
La MAP ordena todos los elementos anteriores en una tabla y distingue entre \textbf{usuarios} (quienes
usan) y \textbf{clientes} (quienes pagan). El problema se escribe con una sintaxis fija:

\begin{center}
\textbf{Usuario} (+ contexto) \conector{necesitan} \textbf{necesidad}
\conector{pero} \textbf{evidencias} \conector{lo que se manifiesta en} \textbf{síntomas}
\conector{porque} \textbf{hallazgos} \conector{[insight]}.
\end{center}

\subsubsection{Procesos de innovación}
\begin{itemize}
  \item \textbf{Design Thinking:} empatizar, definir, idear, prototipar y testear; se puede volver a cualquier etapa.
  \item \textbf{Doble diamante:} descubrir y definir (el problema), desarrollar y entregar (la solución).
        Primero se abre y luego se cierra, dos veces.
\end{itemize}

\subsubsection{Datos}
\begin{itemize}
  \item \textbf{Cualitativos} (entrevistas, observaciones): se \emph{codifican} con etiquetas y luego
        se agrupan por temas (\emph{análisis temático}).
  \item \textbf{Cuantitativos} (encuestas, métricas): media, mediana, moda, varianza y desviación estándar.
  \item Proceso: recolectar, limpiar y organizar, analizar, visualizar, y volver a iterar.
\end{itemize}

\subsubsection{Trabajo en equipo}
Fases de Tuckman (formación, tormenta, normas, desempeño, cierre) y perfiles DISC (dominancia,
influencia, estabilidad, cumplimiento). Hoy Franui es un proyecto de una persona; esto será útil
cuando se sumen más personas.

\subsection{Aplicación a Franui}

\subsubsection{¿Problemática, problema o desafío?}
\begin{itemize}
  \item \textbf{Problemática:} el aprendizaje de idiomas extranjeros en Chile y lo difícil que es mantenerlo en el tiempo.
  \item \textbf{Problema} (el que Franui ataca): personas que quieren aprender idiomas no logran practicar
        vocabulario de forma constante y sienten que no avanzan.
  \item \textbf{Desafío} personal: construir una herramienta propia, bella y útil, y convertirla en empresa.
\end{itemize}

\subsubsection{Persona usuaria (primera segmentación)}
\begin{hipotesis}
\begin{itemize}
  \item \textbf{Edad y ocupación:} estudiantes universitarios y jóvenes profesionales de 18 a 30 años.
  \item \textbf{Lugar:} Chile, en un primer momento; luego, Latinoamérica.
  \item \textbf{Motivo:} estudiar, trabajar o viajar al extranjero, o por gusto personal; muchos quieren más de un idioma.
  \item \textbf{Contexto:} poco tiempo libre; estudian desde el celular, en ratos cortos (micro, filas, pausas).
  \item \textbf{¿Usuario o cliente?} Por definir: la misma persona, su universidad o una empresa podrían pagar.
\end{itemize}
\end{hipotesis}

\subsubsection{Borrador de formulación del problema}
\begin{formulacion}
\textbf{Estudiantes universitarios y jóvenes profesionales en Chile (18–30 años) que quieren aprender uno o más idiomas por su cuenta}
\conector{necesitan} practicar vocabulario y frases de uso cotidiano de forma constante, en ratos cortos y sin miedo a equivocarse,
\conector{pero} \porinvestigar{nivel de inglés de Chile en índices internacionales; porcentaje que abandona apps o cursos de idiomas en las primeras semanas},
\conector{lo que se manifiesta en} \porinvestigar{abandono, sensación de no avanzar, vergüenza al hablar},
\conector{porque} \porinvestigar{causas confirmadas en entrevistas: ejercicios descontextualizados, un solo idioma a la vez, motivación que decae}.
\conector{[Insight]} \porinvestigar{se escribe después de las entrevistas}.
\end{formulacion}

\noindent Este borrador es una hipótesis. Cada \porinvestigar{…} se reemplaza con una fuente (APA 7) o
con un hallazgo de nuestras entrevistas.

\subsubsection{Hallazgos preliminares (hechos)}
Son comentarios de nuestras primeras personas usuarias, en octubre de 2026. Son pocos y vienen de
compañeros y amigos, así que sirven para orientar, no para concluir.

\begin{center}
\small
\begin{tabularx}{\linewidth}{>{\raggedright\arraybackslash}p{.05\linewidth}>{\raggedright\arraybackslash}X>{\raggedright\arraybackslash}X}
\toprule
 & \textbf{Hallazgo} & \textbf{Qué hicimos} \\
\midrule
H1 & 13 personas crearon una cuenta con Google en los dos primeros días. & Corregimos el ranking para que aparezcan todas. \\
H2 & Una persona pidió ejemplos: la traducción sola no basta (\emph{account} = cuenta, ¿del banco o del restaurante?). & Oraciones de ejemplo y notas de Kitty en las palabras confusas. \\
H3 & Una persona valoró la presentación y el progreso, pero sintió demasiada información al inicio. & Inicio más simple para principiantes y tutorial de 2 pasos. \\
H4 & Una persona notó que no hay forma de comprobar si de verdad recuerdas la palabra. & Prueba de Kitty: escribir el significado al final de la sesión. \\
H5 & Varias personas no sabían cómo descargar la app. & Aviso grande de descarga, con pasos para cada celular. \\
H6 & Una persona describió la app como un \enquote{juego dopamínico}, en buen sentido. & Mantener la gamificación. \\
\bottomrule
\end{tabularx}
\end{center}

\subsubsection{Insights preliminares}
\begin{hipotesis}[Insights por validar]
\begin{itemize}
  \item Las personas no solo quieren saber qué significa una palabra, sino \textbf{cuándo usarla}.
  \item El comienzo decide si alguien se queda: \textbf{mucha información al inicio asusta}.
  \item Decir \enquote{lo sé} no basta: \textbf{recordar necesita comprobarse}.
  \item El juego y la recompensa ayudan a volver cada día.
\end{itemize}
\end{hipotesis}

\subsubsection{Sesgos que debemos cuidar}
\begin{itemize}
  \item \textbf{Sesgo de confirmación:} la creadora también es usuaria y quiere que la idea funcione.
  \item \textbf{Muestra cercana:} hasta ahora opinaron amigos y compañeros, que tienden a ser amables.
  \item \textbf{Muestra pequeña:} pocos comentarios; no permiten calcular porcentajes.
\end{itemize}

\subsubsection{Dónde estamos en el doble diamante}
Construimos primero la solución (diamante 2) porque el prototipo nos ayudó a aprender rápido. Ahora
volvemos al diamante 1: \textbf{descubrir y definir} el problema con evidencia, para confirmar o
corregir el rumbo antes de buscar financiamiento.

\subsection{Próximos pasos}
\begin{enumerate}
  \item \textbf{Investigación bibliográfica:} nivel de idiomas en Chile, abandono en apps de idiomas,
        tamaño del mercado de aprendizaje en línea. Cada dato con su fuente.
  \item \textbf{Observación:} mirar a 3–5 personas mientras estudian un idioma (con app, clases o por su cuenta) y tomar notas.
  \item \textbf{Entrevistas de validación:} 5–8 personas fuera de nuestro círculo cercano, con un guion preparado.
  \item \textbf{Encuesta corta} a estudiantes, para tener datos cuantitativos.
  \item \textbf{Análisis temático} de todo el feedback: codificar y agrupar por temas.
  \item \textbf{Completar la MAP}, separando usuarios y clientes.
  \item \textbf{Medir el uso real} de forma anónima: personas que vuelven al día 1, al día 7 y al día 30.
\end{enumerate}


% ------------------------- Registro de decisiones -------------------------
% Registro de decisiones: una fila por decisión importante (la más nueva arriba).
\clearpage
\section{Registro de decisiones}

\begin{center}
\small
\begin{tabularx}{\linewidth}{>{\raggedright\arraybackslash}p{.15\linewidth}>{\raggedright\arraybackslash}X>{\raggedright\arraybackslash}X}
\toprule
\textbf{Fecha} & \textbf{Decisión} & \textbf{Por qué} \\
\midrule
10-nov-2026 & Usar los conceptos del curso Proyectos como guía para formalizar Franui, sin que sea una tarea. & Las herramientas (MAP, doble diamante, validación) sirven para construir la startup y postular a fondos. \\
10-nov-2026 & Llevar esta bitácora en LaTeX, dentro del repositorio del proyecto. & Queda guardada con su historial de cambios y se puede compartir como PDF. \\
31-oct-2026 & Estética minimalista para la app y la web, con la misma paleta. & Referencias de marcas premium; más elegante y fácil de leer. \\
31-oct-2026 & Quitar del banco las expresiones universales (\emph{oh, oh dear, wow, hooray}). & No se aprenden: se entienden en cualquier idioma. \\
14-oct-2026 & Mudar la app a \url{https://polykat-franui.github.io}. & Una dirección con el nombre del proyecto, sin perder el avance de nadie. \\
14-oct-2026 & Toda traducción generada por IA queda como borrador hasta revisarla. & La calidad del contenido es la base de la confianza. \\
\bottomrule
\end{tabularx}
\end{center}


\end{document}
